\section{Limitations}\label{sec:limits}

Every run uses one seed, and R@1 on RSTPReid carries a standard error of about 1.1 points, so differences below roughly two points (LoRA versus full fine-tuning with SDM) should be read as parity; the SDM-versus-InfoNCE and SDM-versus-MPCE gaps are about three standard errors and are mirrored in the less noisy mAP. The SDM run used a shorter learning-rate schedule than the other two; \cref{sec:ablation} argues that this does not explain the result, but a matched rerun would remove the caveat. All conclusions concern RSTPReid with random batch sampling. On larger datasets, or with an identity-balanced sampler that places several images of each person in a batch, the multi-positive targets of SDM and MPCE become frequent and may then matter. The IRRA numbers are quoted from the paper rather than reproduced, and the mechanism proposed in \cref{sec:ablation} has not yet been isolated experimentally. Finally, no federated experiment is reported: \cref{sec:fl} is an accounting argument, and how accuracy behaves under non-IID clients remains to be measured in the thesis.

\section{Conclusion}\label{sec:conclusion}

This report tested, in the centralised setting, the premise behind using PEFT for federated text-based person search: that a far smaller trainable part loses little accuracy. Starting from IRRA, we asked how much accuracy a frozen CLIP backbone retains with LoRA and the SDM loss alone, and why SDM helps. With \trainPct\% of the parameters, \sizeRatio$\times$ fewer than full fine-tuning, LoRA+SDM matches IRRA's fully fine-tuned SDM-only model on RSTPReid (\sdmPeakR\% peak, \sdmConvR\% converged, against 57.20\%) and trails the complete IRRA model by \gapIrraR{} points of R@1, while cutting the per-round upload of a \kClients-client federation from \upFullRoundGiB\,GiB to \upLoraRoundMiB\,MiB. The premise holds. On the second question our evidence contradicts the usual explanation. Under random sampling SDM's targets are one-hot in \oneHotFrac\% of rows; in that regime SDM is a scaled mean-absolute-error loss plus an entropy bonus rather than InfoNCE, and a control loss with SDM's targets but the forward KL behaves exactly like InfoNCE. The benefit of SDM therefore lies in the direction of its divergence, which resists the overfitting that InfoNCE shows on this small dataset. Next steps are seeded reruns with a common schedule, ablations of the two terms of \cref{eq:decomp}, an identity-balanced sampler, and, in the thesis, federated training of LoRA+SDM over camera-level clients on RSTPReid, CUHK-PEDES and ICFG-PEDES.
